% Insert after Appendix C (EM view of binary success regularization)
% and before Appendix D (Training and implementation details).
% The main-text pointer at Eq.~(7) is Appendix~\ref{app:critic-gate}.
\section{Estimating the critic gate}
\label{app:critic-gate}

The critic gate in Eq.~(7) compares two local sensitivities of the
existing contrastive score $f(s,a,g)$. Both quantities are estimated
from a small number of additional critic evaluations at the observed
pair $(s,g)$ drawn from $\mathcal{D}_{\mathrm{succ}}$. The stored
successful action is not among the probes.

Action sensitivity $\sigma_a(s,g)$ is the standard deviation of
$f(s,a',g)$ over eight actions $a'$ drawn independently and uniformly
from the environment action space, with $(s,g)$ held fixed.

State sensitivity $\sigma_s(s,g)$ is the standard deviation of the
score under a local change of $s$ alone. We draw eight perturbations
\begin{equation*}
\tilde{s}
=
s + \varepsilon\,\hat{\sigma}\odot z,
\qquad
z\sim\mathcal{N}(0,I),
\end{equation*}
evaluate $f(\tilde{s},a_{\pi},g)$ with the same goal and a single
policy action $a_{\pi}\sim\pi(\cdot\mid s,g)$, and take the standard
deviation of those scores. The vector $\hat{\sigma}$ is the
coordinate-wise standard deviation of states in the current success
minibatch and is treated as a constant; we use $\varepsilon=0.1$.
Because the perturbation is local to the queried state, $\sigma_s$
measures how the score changes when that state is nudged, rather than
how scores vary across different states in the minibatch.

The resulting weight $w(s,g)$ multiplies the success likelihood in
Eq.~(8). Gradients are stopped through $w$. If the two standard
deviations vanish together, the gate is set to one.
